% ----------------------------------------------------------------------------
%
% ----------------------------------------------------------------------------

\documentclass[11pt, parskip=half]{scrartcl}       % KOMA-Skript für Artikel

%% Präambel
\usepackage[english, ngerman]{babel} % deutsche typogr. Regeln + Trenntabelle
\usepackage[T1]{fontenc}             % interner TeX-Font-Codierung
\usepackage{lmodern}                 % Font Latin Modern
\usepackage[utf8]{inputenc}          % Font-Codierung der Eingabedatei
\usepackage[babel]{csquotes}         % Anführungszeichen
\usepackage{graphicx}                % Graphiken
\usepackage{booktabs}                % Tabellen schöner
\usepackage{amsmath}      % Mathematik
\usepackage{amssymb}      % Mathematische Symbole
\usepackage{float}
\usepackage[pdftex]{hyperref}
\usepackage{subcaption}
\usepackage{url}
\hypersetup{
  bookmarksopen=true,
  bookmarksopenlevel=3,
  colorlinks,
  citecolor=blue,
  linkcolor=blue
}
\usepackage{scrhack} % unterdrückt Fehlermeldung von listings
\usepackage[backend=bibtex]{biblatex}
\addbibresource{/home/niklas/Documents/bibfile/bibliographie.bib}


%% Nummerierungstiefen
\setcounter{tocdepth}{3}     % 3 Stufen im Inhaltsverzeichnis
\setcounter{secnumdepth}{3}  % 3 Stufen in Abschnittnummerierung

\usepackage[nolist]{acronym}
\begin{acronym}
\acro{pwa}[PWA]{Progressive Web App}
\end{acronym}

\newcommand\litem[1]{\item{\bfseries#1.\space}}

\newcommand*{\N}{\mathbb{N}}
\newcommand*{\Z}{\mathbb{Z}}
\newcommand*{\Q}{\mathbb{Q}}
\newcommand*{\R}{\mathbb{R}}

% Code packages
\usepackage{listingsutf8} % Code mit UTF-8 Support
\usepackage{color,xcolor} % Farben definierbar als HTML, RGB, ...
\usepackage{textcomp}

\definecolor{alabasterGray}{HTML}{F7F7F7}
\definecolor{alabasterBlue}{HTML}{325CC0}
\definecolor{alabasterGreen}{HTML}{448C27}
\definecolor{alabasterPink}{HTML}{7A3E9D}
\definecolor{alabasterRed}{HTML}{AA3731}

\lstset{
  basewidth={0.5em,0.45em},
  extendedchars=true,
  backgroundcolor=\color{alabasterGray}, % background color (color or xcolor needed)
  basicstyle=\small\ttfamily,
  keywordstyle=\color{alabasterBlue},
  commentstyle=\color{alabasterRed},
  rulecolor=\color{black},          % frame color may change if not set (frame on comment, is comment-colored)
  stringstyle=\color{alabasterGreen},
  numberstyle=\tiny\color{gray},    %
  numbers=none,                     % where to put the line-numbers (none, left, right)
  numbersep=7pt,                    % margin between numbers and code
  stepnumber=1,                     % every n-th row will be numbered
  captionpos=b,                     % sets the caption-position to bottom
  frame=single,                     % adds a frame around the code
  keepspaces=true,                  % keeps spaces in text
  showtabs=false,                   % show tabs as special char
  showspaces=false,                 % show spaces as special char
  showstringspaces=false,           % show spaces in strings only as special char
  tabsize=2,                        % tabsize is 2 spaces
  breakatwhitespace=false,          % sets if automatic breaks should only happen at whitespace
  breaklines=true,                  % sets automatic line breaking
}
\lstset{
  literate={ö}{{\"o}}1
           {ä}{{\"a}}1
           {ü}{{\"u}}1
           % https://tex.stackexchange.com/questions/17739/listings-package-how-to-highlight-math-operators
           {true}{{{\color{alabasterPink}true}}}4
           {false}{{{\color{alabasterPink}false}}}5
           {TRUE}{{{\color{alabasterPink}TRUE}}}4
           {FALSE}{{{\color{alabasterPink}FALSE}}}5
           {True}{{{\color{alabasterPink}True}}}4
           {False}{{{\color{alabasterPink}False}}}5
}
\expandafter\def\expandafter\UrlBreaks\expandafter{\UrlBreaks\do\-\do\\/}

\begin{document}

\titlehead{\includegraphics[width=\textwidth]{src/Logo_THM_CG_FB06.png}}

%% Titelseite
\subject{Manusskript}
\title{Progressive Web Apps}
\subtitle{Kurs \enquote{Hauptseminar -- Mobile Techonologies}}
\author{Niklas Deworetzki}
\date{\today}
\maketitle

\vspace*{1.5cm}
\section*{\centerline{Zusammenfassung}}

Diese Arbeit befasst sich mit Progressive Web Applications.


\newpage

\tableofcontents
\newpage

\section{Einleitung}

% TODO Einleitung
https://medium.com/@amberleyjohanna/seriously-though-what-is-a-progressive-web-app-56130600a093

\acp{pwa} werden immer relevanter.
Ergänzendes Element zwischen herkömmlicher Anwendung und Website.
Was kann man damit machen? Worum geht's eigentlich?

\section{Was ist eine \ac{pwa}?}

Der Begriff \acl{pwa} beschreibt eine Kategorie von Webanwendungen, welche eine Reihe an Eigenschaften erfüllen, um das Erlebnis einer nativen Anwendung in modernen Browsern zu bieten.
Geprägt wurde dieser Begriff von Alex Russell, welcher in seinem Blog\cite{russell_pwaescapingtabs} beschreibt, wie durch die Weiterentwicklung und Standardisierung von Browsern und Webtechnologien eine neue Art von Anwendung entstanden ist.
Diese Kategorie von Anwendungen entspreche dem nächsten Schritt in einer Reihe von Technologien, die versuchen, Entwicklern die Möglichkeit zu geben, mit nur einer plattformübergreifenden Anwendung auf plattformspezifische Ressourcen zuzugreifen.

Diese Kategorie von Webanwendungen ist also eine Art Nachfolger auf Technologien wie \enquote{Universal Windows Platform}-Apps\cite{msdocs_uwp} von Microsoft oder die von W3C standardisierten \enquote{Packaged Web Apps}\cite{w3c_packagedwebapps}, welche auch als \enquote{Widget} bekannt sind.
Der Begriff \enquote{Progressive} in \acl{pwa} steht für \enquote{Progressive Enhancement} (fortschreitende Verbesserung) und beschreibt ebendiese Vorgehensweise, eine Webanwendung auf möglichst vielen Geräten mit verschiedenen Kapazitäten nutzbar zu machen.
Die Idee dabei ist es, eine Grundversion der Anwendung zu erstellen, die auf jedem Gerät funktionieren kann.
Weitere Funktionalität für potentere Geräte werden dann auf diese Grundversion aufgebaut.

Da der Begriff der \ac{pwa} zusammen mit den notwendigen Technologien über einen längeren Zeitraum gewachsen ist, gibt es keine feste Definition für das, was eine \ac{pwa} ausmacht.
Die Beschreibung einer \ac{pwa} basiert vielmehr auf einer Menge an Eigenschaften, die eine Anwendung innehaben kann.
Im Folgenden werden daher einige Eigenschaften zusammengetragen und erläutert, die immer wieder für die Beschreibung von \acp{pwa} verwendet werden.
Ein besonderer Schwerpunkt wird dabei auf die Standpunkte der Google- und Mozilla-Entwickler gelegt, da diese treibende Kraft beim Entwickeln von \acp{pwa} sind und gemeinsam die Fähigkeiten dieser Gruppe von Anwendungen ausbauen.

\newpage % TODO: Eine Seite mehr an dieser Stelle. Nachher vielleicht anpassen?

Die gesammelten Eigenschaften einer \ac{pwa} lassen sich dabei in folgende Kategorien einteilen, in welcher sie auch erläutert werden.

\begin{enumerate}
  \litem{Grundlegende Eigenschaften} Eigenschaften einer \ac{pwa}, die dem Charakter einer \ac{pwa} entsprechen und ohne die ein Funktionieren der Anwendung nicht möglich wäre.

  \litem{Verbreitete Eigenschaften} Eigenschaften einer \ac{pwa}, welche von verschiedenen Quellen genannt werden, aber nicht unbedingt notwendig für das Funktionieren der Anwendung sind.

  \litem{Designmerkmale} Eigenschaften einer \ac{pwa}, die sich positiv auf das Nutzererlebnis auswirken.
\end{enumerate}


\subsection{Grundlegende Eigenschaften}

Die wohl wichtigste Eigenschaft einer \ac{pwa} ist, dass sie installierbar sein muss.
Wenn ein Nutzer eine Website besucht, welche als \ac{pwa} agiert, so kann dem Nutzer eine Installationsaufforderung gezeigt werden.
Durch den Installationsvorgang werden teile der Anwendung auf dem lokalen Gerät gespeichert und die \ac{pwa} optisch zu den nativen Anwendung integriert.
Ob diese Integration nun über das Hinzufügen eines Desktopicons, eine Verknüpfung auf dem Startbildschirm oder auch nur über einen Eintrag in einem Anwendungslauncher erfolgt, ist dabei vom unterliegenden System abhängig.
Auf einem Android-Gerät wird es beispielsweise sinnvoll sein, die \ac{pwa} als Icon auf dem Startbildschirm anzuzeigen.
Auf einem Windows-Gerät stattdessen wird dem Nutzer eine Verknüpfung auf dem Desktop lieber sein.
Da die \ac{pwa} nicht eigenständig entscheiden kann, welches Gerät der Nutzer verwendet und wohl unmöglich die Fähigkeit besitzt, für jedes Gerät die entsprechende Konfiguration durchzuführen, muss der Installationsvorgang durch den Browser unterstützt und verwaltet werden.

Beim anschließenden Ausführen einer installierten \ac{pwa} ist erneut Unterstützung vom Browser notwendig.
\acp{pwa} als Webanwendung werden mit Webtechnologien wie HTML und CSS aufgebaut und erhalten ihre Funktionalität durch JavaScript.
Betriebsysteme bieten jedoch nur eine begrenzte Unterstützung für Programmformate, die nativ ausgeführt werden können.
Abgesehen von Skriptdateien wie etwa Shell-Scripts für Linux oder Command-Files für Windows, die in einem betriebsystemabhängigen Format geschrieben sein müssen, wird meist nur die Ausführung von Binärformaten unterstützt\cite{fisher_executablelist}.
Es wird also Hilfe vom Browser benötigt, welcher die installierte \ac{pwa} verwaltet und daher alle benötigten Dateien laden und ausführen kann.
Zudem ist es dem Browser möglich, für die Ausführung einer \ac{pwa} bestimmte Regeln festzulegen.
In Google-Chrome beispielsweise wird jede \ac{pwa} in einem eigenen Fenster gestartet, wodurch das Nutzererlebnis einer nativen Anwendung entstehen soll\cite{googledevs_pwa}.


Durch den Installationsprozess werden nur Teile \ac{pwa} lokal auf dem Gerät gespeichert.
Der Grund dafür ist, dass es meist nicht genügt, eine Website einfach vollständig zu kopieren, um sie offline verfügbar zu machen.
Viele Webanwendungen, wie beispielsweise die \ac{pwa} von Twitter\cite{twitter_pwa}, basieren darauf stetig neue Inhalte von einem Server abzufragen.
Die Verwaltung, welche Inhalte offline verfügbar sein sollen und welche nicht, geschieht über sogenannte \enquote{Service Worker}.
Dadurch dass jede \ac{pwa} lokale und entfernte Inhalte verwalten muss, ist das vorhandensein von Service Workern, eine weitere Eigenschaft, die jede \ac{pwa} besitzt.

Eine Voraussetzung der Service Worker erzwingt eine weitere Eigenschaft aller \acp{pwa}.
Service Worker können laut Spezifikation nur auf Seiten ausgeführt werden, die über HTTPS ausgeliefert werden.
Der Grund dafür ist, dass Service Worker die Fähigkeit besitzen, Anfragen zu filtern, umzuleiten und auszulesen.
Die verschlüsselte Verbindung bei HTTPS verhindert dabei die Manipulation durch Dritte beim Ausliefern des Service Worker.
So soll sichergestellt werden, dass Service Worker nicht für das Ausspionieren von Daten genutzt werden\cite{ServiceWorker_explained}.

% TODO: Abschnitt URLs
Die letzte grundlegende Eigenschaft einer \ac{pwa} ergibt sich aus der Herkunft dieser.
Wie bei einer herkömmlichen Website geschieht die Auswahl der angezeigten Inhalte über URLs.
Daraus ergibt sich, dass jeder Teil der Anwendung über eine URL erreichbar sein muss.


\subsection{Häufige Eigenschaften}



\subsection{Designmerkmale}


\section{Technische Voraussetzungen}

\subsection{Installationsformat}

\subsection{Datenverwaltung durch Service Worker}

\subsection{Interaktion mit Nutzern}

\subsection{Interaktion mit Hardware}

\subsection{Browserspezifische Voraussetzungen}


\section{Vorteile von PWAs}

\section{Kritik an PWAs}


\section{Fazit}

\newpage

\printbibliography

\end{document}

%%% Local Variables:
%%% mode: latex
%%% TeX-master: t
%%% End:
